\uuid{mvpS}
\titre{Angle}
\auteur{bodin ; hanani ; mzari}
\organisation{exo7}
\licence{CC BY-NC-SA 4.0}
\datecreate{2020-01-20}
\niveau{L1}
\module{Géométrie}
\chapitre{Géométrie du plan}
\sousChapitre{Distance, norme, produit scalaire}
\difficulte{2}
\theme{angle}
\exo7qcmid{lille-1.geometrie-plane.011}

\contenu{
\begin{question}
\qclassification{140.01,140.02}
\qtags{motcle=angle}
Dans le plan muni d'un repère orthonormé direct $(O,\vec{i},\vec{j})$, on considère les vecteurs $\displaystyle \vec{u}=\left(1,1\right)$ et $\displaystyle \vec{v}=\left(1,\sqrt{3}\right)$. Quel est la mesure $\alpha \in [0,2\pi[$ de l'angle orienté entre $\vec{u}$ et $\vec{v}$ ?
\begin{answers}
    \bad{$\displaystyle \alpha =\frac{\pi}{4}$}
    \bad{$\displaystyle \alpha =\frac{\pi}{3}$}
    \good{$\displaystyle \alpha =\frac{\pi}{12}$}
    \bad{$\displaystyle \alpha =\frac{7\pi}{12}$}
\end{answers}
\begin{explanations}
Une mesure de l'angle orienté entre $\vec{i}$ et $\vec{u}$ est $\displaystyle a=\frac{\pi}{4}$ et une mesure de l'angle orienté entre $\vec{i}$ et $\vec{v}$ est $\displaystyle b=\frac{\pi}{3}$. Donc $\displaystyle \alpha =b-a=\frac{\pi}{12}$.
\end{explanations}
\end{question}
}
